% TOP_PROBLEM: {"id": 20000957, "title": "Four concyclic points and bounded Euclidean Ramsey witnesses", "category_id": 2, "category": {"id": 2, "name": "combinatorics", "display_name": "Combinatorics", "description": "Counting problems, graph theory, discrete structures.", "slug": "combinatorics", "order_index": 2, "created_at": "2026-07-31T15:26:25.671Z"}, "status": "partially_solved"}
% FIRST_POSTED: 2026-10-06
% ATTEMPT_STATUS: solved
% OPENAI_CLAIM: {"schema_version":1,"source_repo":"https://github.com/openai/math","source_commit":"adc7f1241b42e322a6451854ab7e4b4c146bf78a","release_date":"2026-10-06","match":"stronger","resolution":"proved","independently_reviewed":false,"second_pass":"agreed","families":[{"family":"172","title":"Classification of finite Euclidean Ramsey configurations","subject":"Combinatorics","manuscripts":[{"title":"A classification of finite Euclidean Ramsey configurations","pdf_path":"preprints/A-classification-of-finite-Euclidean-Ramsey-configurations-September-23-2026/paper.pdf","readme_path":"preprints/A-classification-of-finite-Euclidean-Ramsey-configurations-September-23-2026/README.md","date":"2026-09-23","theorem_ref":"thm:classification (sections/01-introduction.tex:74); cons:kites (sections/07-consequences.tex:124); cons:quadrilateral (sections/07-consequences.tex:106)","statement_sources":[{"path":"preprints/A-classification-of-finite-Euclidean-Ramsey-configurations-September-23-2026/build/sections/01-introduction.tex","line":74,"label":"thm:classification"},{"path":"preprints/A-classification-of-finite-Euclidean-Ramsey-configurations-September-23-2026/build/sections/07-consequences.tex","line":124,"label":"cons:kites"},{"path":"preprints/A-classification-of-finite-Euclidean-Ramsey-configurations-September-23-2026/build/sections/07-consequences.tex","line":106,"label":"cons:quadrilateral"}]}],"lean":{"doc_path":"lean/docs/172.md","comparators":["lean/ComparatorChallenges/EuclideanRamsey.json","lean/ComparatorChallenges/EuclideanRamsey.lean","lean/ComparatorChallenges/EuclideanRamseyCircle.json","lean/ComparatorChallenges/EuclideanRamseyCircle.lean","lean/ComparatorChallenges/EuclideanRamseyNine.json","lean/ComparatorChallenges/EuclideanRamseyNine.lean","lean/ComparatorChallenges/EuclideanRamseyQuadratic.json","lean/ComparatorChallenges/EuclideanRamseyQuadratic.lean","lean/ComparatorChallenges/EuclideanRamseySpherical.json","lean/ComparatorChallenges/EuclideanRamseySpherical.lean","lean/ComparatorChallenges/EuclideanRamseyTransitive.json","lean/ComparatorChallenges/EuclideanRamseyTransitive.lean","lean/ComparatorChallenges/GrahamSpherical.json","lean/ComparatorChallenges/GrahamSpherical.lean"],"declarations":["OAI.EuclideanRamsey.classification","OAI.GrahamSpherical.full_main"],"covers_main_theorem":true},"reasoning_trace":null}]}
\documentclass[11pt]{article}
\usepackage[margin=25mm]{geometry}
\usepackage{fontspec}
\setmainfont{Latin Modern Roman}
\usepackage{amsmath,amssymb,mathtools}
\usepackage{unicode-math}
\setmathfont{Latin Modern Math}
\usepackage{enumitem}
\usepackage{xurl,hyperref,bookmark}
\hypersetup{colorlinks=true,urlcolor=blue}
\setcounter{secnumdepth}{0}
\setlength{\parindent}{0pt}
\setlength{\parskip}{5pt}
\setlength{\emergencystretch}{3em}
\begin{document}
\section*{Four concyclic points and bounded Euclidean Ramsey witnesses}

\subsection{OpenAI's claim}
\textit{Finite Euclidean Ramsey classification}: The announced characterization of finite Euclidean Ramsey sets is a tensor condition over the coordinate field, at the original congruence scale. Consequences include Ramsey behavior for subtransitive sets and small circular sets, as well as circular Ramsey examples that are not subtransitive.

\subsection{Mathematical statement of the claimed result}
\paragraph{Mathematical claim 1.} For distinct points $A=\{a_1,\ldots,a_s\}\subset\mathbb R^d$ with $s\ge2$ and affine span $\mathbb R^d$, put $p_i=(1,a_i)^{\mathsf T}$, $F=\mathbb Q((a_i)_\alpha)$, $B=F\otimes_{\mathbb Q}F$, and $m_F(x\otimes y)=xy$. The Ramsey quantifiers are $\forall r\in\mathbb Z_{\ge2}\ \exists D\in\mathbb Z_{\ge1}\ \forall c:\mathbb R^D\to\{1,\ldots,r\}$: some isometric embedding $\iota:A\to\mathbb R^D$ has $c\circ\iota$ constant. No measurability of $c$ or rescaling of $A$ is permitted. The classification claims this is equivalent to existence of $P\in\operatorname{Mat}_{d+1}(B)$ satisfying \[(p_i\otimes1)^{\mathsf T}P(1\otimes p_i)=0\quad(1\le i\le s),\qquad m_F(P_{\alpha\beta})=\delta_{\alpha\beta}\quad(1\le\alpha,\beta\le d).\] The constant-coordinate row and column have no additional multiplied-entry conditions. Corollaries say every set of at most five circle points is Ramsey, and for transcendental $a\in(-1,1)$ the cyclic kite $\{(-1,0),(1,0),(a,\sqrt{1-a^2}),(a,-\sqrt{1-a^2})\}$ is Ramsey while embedding in no finite transitive Euclidean set.

\subsection{Outline of the argument}
\textit{Finite Euclidean Ramsey classification}: Tensor cancellation is first realized by exact pairs of permutation constructions. A group-endpoint criterion and finite-string realization prepare the distinct weights needed for a Hales--Jewett line to have precisely unit distance scale. Further averaging and short corrective paths equalize endpoints, making the sufficient tensor condition usable under arbitrary colorings.

\subsection{Scope relative to this problem}
\textit{Finite Euclidean Ramsey classification}: The corollary establishes that every set of at most five concyclic points is Ramsey, which includes every four-point subset in the explicit recorded conjecture; the trailing rival-conjecture fragment contains no additional mathematical clause.

The claimed result implies the stated question and also supplies the additional strength described above.

\subsection{Formal verification}
Family 172 describes Lean coverage for the tensor criterion and the stated spherical, subtransitive, and circular configuration consequences.

\textit{Finite Euclidean Ramsey classification}: The manuscript links \texttt{EuclideanRamsey.lean} (declaration \texttt{OAI.EuclideanRamsey.classification}), \texttt{EuclideanRamseyCircle.lean}, \texttt{EuclideanRamseyNine.lean}, \texttt{EuclideanRamseyQuadratic.lean}, \texttt{EuclideanRamseySpherical.lean}, \texttt{EuclideanRamseyTransitive.lean}, \texttt{GrahamSpherical.lean} (declaration \texttt{OAI.GrahamSpherical.full\_main}).

The documented selection includes this manuscript’s main assertion; that assertion retains the restrictions stated above.

This site has not run the Lean development. These links describe OpenAI’s selected formal statements and should be read with their stated scope.

\subsection{Status caveat}
OpenAI reports that an internal model produced the manuscripts. Their arguments have not been independently verified as proofs by this site. Any portion outside the documented formal scope may contain mistakes, and the catalogue label reports the announced mathematical claim.

\subsection{Sources}
\begin{itemize}
\item OpenAI, \emph{A classification of finite Euclidean Ramsey configurations}, 2026-09-23; citation key \texttt{OAI:A-classification-of-finite-Euclidean-Ramsey-configurations-September-23-2026}. \href{https://github.com/openai/math/blob/adc7f1241b42e322a6451854ab7e4b4c146bf78a/preprints/A-classification-of-finite-Euclidean-Ramsey-configurations-September-23-2026/paper.pdf}{PDF (GitHub)}; \href{https://github.com/openai/math/raw/adc7f1241b42e322a6451854ab7e4b4c146bf78a/preprints/A-classification-of-finite-Euclidean-Ramsey-configurations-September-23-2026/paper.pdf}{PDF (direct)}; \href{https://github.com/openai/math/blob/adc7f1241b42e322a6451854ab7e4b4c146bf78a/preprints/A-classification-of-finite-Euclidean-Ramsey-configurations-September-23-2026/README.md}{Manuscript page and citation}; \href{https://github.com/openai/math/blob/main/preprints/A-classification-of-finite-Euclidean-Ramsey-configurations-September-23-2026/paper.pdf}{Latest version}; Statement source: \href{https://github.com/openai/math/blob/adc7f1241b42e322a6451854ab7e4b4c146bf78a/preprints/A-classification-of-finite-Euclidean-Ramsey-configurations-September-23-2026/build/sections/01-introduction.tex\#L74}{thm:classification}; Statement source: \href{https://github.com/openai/math/blob/adc7f1241b42e322a6451854ab7e4b4c146bf78a/preprints/A-classification-of-finite-Euclidean-Ramsey-configurations-September-23-2026/build/sections/07-consequences.tex\#L124}{cons:kites}; Statement source: \href{https://github.com/openai/math/blob/adc7f1241b42e322a6451854ab7e4b4c146bf78a/preprints/A-classification-of-finite-Euclidean-Ramsey-configurations-September-23-2026/build/sections/07-consequences.tex\#L106}{cons:quadrilateral}.
\item Family 172: \href{https://github.com/openai/math/blob/adc7f1241b42e322a6451854ab7e4b4c146bf78a/lean/docs/172.md}{Lean scope}.
\item \href{https://github.com/openai/math/blob/adc7f1241b42e322a6451854ab7e4b4c146bf78a/lean/ComparatorChallenges/EuclideanRamsey.json}{Comparator EuclideanRamsey.json}.
\item \href{https://github.com/openai/math/blob/adc7f1241b42e322a6451854ab7e4b4c146bf78a/lean/ComparatorChallenges/EuclideanRamsey.lean}{Comparator EuclideanRamsey.lean}.
\item \href{https://github.com/openai/math/blob/adc7f1241b42e322a6451854ab7e4b4c146bf78a/lean/ComparatorChallenges/EuclideanRamseyCircle.json}{Comparator EuclideanRamseyCircle.json}.
\item \href{https://github.com/openai/math/blob/adc7f1241b42e322a6451854ab7e4b4c146bf78a/lean/ComparatorChallenges/EuclideanRamseyCircle.lean}{Comparator EuclideanRamseyCircle.lean}.
\item \href{https://github.com/openai/math/blob/adc7f1241b42e322a6451854ab7e4b4c146bf78a/lean/ComparatorChallenges/EuclideanRamseyNine.json}{Comparator EuclideanRamseyNine.json}.
\item \href{https://github.com/openai/math/blob/adc7f1241b42e322a6451854ab7e4b4c146bf78a/lean/ComparatorChallenges/EuclideanRamseyNine.lean}{Comparator EuclideanRamseyNine.lean}.
\item \href{https://github.com/openai/math/blob/adc7f1241b42e322a6451854ab7e4b4c146bf78a/lean/ComparatorChallenges/EuclideanRamseyQuadratic.json}{Comparator EuclideanRamseyQuadratic.json}.
\item \href{https://github.com/openai/math/blob/adc7f1241b42e322a6451854ab7e4b4c146bf78a/lean/ComparatorChallenges/EuclideanRamseyQuadratic.lean}{Comparator EuclideanRamseyQuadratic.lean}.
\item \href{https://github.com/openai/math/blob/adc7f1241b42e322a6451854ab7e4b4c146bf78a/lean/ComparatorChallenges/EuclideanRamseySpherical.json}{Comparator EuclideanRamseySpherical.json}.
\item \href{https://github.com/openai/math/blob/adc7f1241b42e322a6451854ab7e4b4c146bf78a/lean/ComparatorChallenges/EuclideanRamseySpherical.lean}{Comparator EuclideanRamseySpherical.lean}.
\item \href{https://github.com/openai/math/blob/adc7f1241b42e322a6451854ab7e4b4c146bf78a/lean/ComparatorChallenges/EuclideanRamseyTransitive.json}{Comparator EuclideanRamseyTransitive.json}.
\item \href{https://github.com/openai/math/blob/adc7f1241b42e322a6451854ab7e4b4c146bf78a/lean/ComparatorChallenges/EuclideanRamseyTransitive.lean}{Comparator EuclideanRamseyTransitive.lean}.
\item \href{https://github.com/openai/math/blob/adc7f1241b42e322a6451854ab7e4b4c146bf78a/lean/ComparatorChallenges/GrahamSpherical.json}{Comparator GrahamSpherical.json}.
\item \href{https://github.com/openai/math/blob/adc7f1241b42e322a6451854ab7e4b4c146bf78a/lean/ComparatorChallenges/GrahamSpherical.lean}{Comparator GrahamSpherical.lean}.
\item \href{https://github.com/openai/math/blob/adc7f1241b42e322a6451854ab7e4b4c146bf78a/CONTENTS.md}{OpenAI catalogue entry}.
\end{itemize}
\end{document}
